\documentclass{BeamerTemplate}
\title{Chapitre 5 - Modèles linéaires généralisés mixtes (MLGM)}
\subtitle{STT-4300}
\date{\today}

\IfFileExists{../STT-4300-references.bib}{%
	\addbibresource{../STT-4300-references.bib}%
}{%
	\addbibresource{STT-4300/STT-4300-references.bib}%
}

\begin{document}
\begin{refsection}

\begin{frame}[t]
\titlepage
\end{frame}

\section{Modèles linéaires généralisés mixtes (MLGM)}
\label{GLMM.section}

\subsection{Définition du modèle}
\begin{frame}{Modèles linéaires généralisés mixtes (MLGM)}
Avec des données longitudinales ou groupées, il faut tenir compte de la
{\color{blue}corrélation intra-grappe} (\textit{within-cluster correlation}).

\bigskip

Le MLG n'est pas adéquat. Le {\color{blue}MLGM} (\textit{generalized linear mixed model}, GLMM) étend le MLG en permettant aux coefficients de varier selon la grappe, ce qui tient compte de la corrélation intra-grappe.

\bigskip
Le MLGM est un compromis entre le modèle de {\color{blue}moyenne populationnelle} (\textit{population-average model}) du MLG et une approche de modélisation (sur-pa\-ra\-mé\-tri\-sée) qui ajusterait chaque grappe séparément.
\end{frame}

\begin{frame}
Ingrédients :
\begin{itemize}
\item vecteurs de réponse $Y_i = (Y_{i1},\ldots,Y_{in_i})$, pour $i=1,\ldots,n$.
\item matrice de design (\textit{design matrix}) $X_i$ pour les {\color{blue}effets fixes} (\textit{fixed effects})
\item matrice de design $Z_i$ pour les {\color{blue}effets aléatoires} (\textit{random effects}) (pouvant chevaucher
$X_i$).
\item effets aléatoires $\gamma_i$ pour chaque grappe $i$. Ce sont des variables aléatoires {\color{blue}latentes} (\textit{latent}) (non observées).
\end{itemize}

\bigskip
Hypothèses :
\begin{itemize}
\item $Y_{it} \mid \gamma_i$ sont indépendamment distribués selon une
  distribution de la {\color{blue}famille exponentielle} (\textit{exponential family}).
\item $\gamma_i$ sont indépendamment distribués selon ${\mathcal
    N}(0, \Psi)$.
\end{itemize}

\bigskip
Deux sources d'erreur : le processus d'échantillonnage et le processus de mesure
au sein de l'échantillon lui-même.
\end{frame}

\begin{frame}{Cas gaussien avec lien identité}
Dans le cas gaussien avec {\color{blue}lien identité} (\textit{identity link}), le modèle peut s'écrire :
$$ Y_i = X_i \beta + Z_i \gamma_i + \epsilon_i \ \ \  (\mbox{ou } Y_{it} = x_{it}^T \beta + z_{it}^T \gamma_i + \epsilon_{it}),$$
avec $\epsilon_i \sim {\mathcal N}(0, \sigma^2_{\epsilon} I)$ et $\gamma_i  \sim
{\mathcal N}(0,\Psi)$, par exemple $\Psi = \sigma^2_{\gamma} I$.

\bigskip
De manière équivalente : $E(Y_i \mid \gamma_i) = X_i \beta + Z_i \gamma_i$.


\bigskip

Il s'ensuit que
$$Y_i \sim {\mathcal N}(X_i \beta, Z_i \Psi Z_i^T + \sigma^2_{\epsilon} I),$$
et la (log-)vraisemblance peut être écrite explicitement pour cette
distribution normale multivariée (mais nécessite tout de même un algorithme
d'optimisation en raison des paramètres dans $\Psi$).
\end{frame}

\begin{frame}

\begin{center}
\begin{tabular}{cc} 
Ordonnée à l'origine aléatoire (\textit{random intercept}) & Pente aléatoire (\textit{random slope}) \\
$Y_{it} = x_{it}^T \beta + \gamma_i + \epsilon_{it}$ & $Y_{it} = x_{it}^T \beta + z_{it}^T \gamma_i + \epsilon_{it}$\\
\includegraphics[width = 5.5cm]{figures/RandomInterceptGaussian.pdf}  & \includegraphics[width =5.5cm]{figures/RandomSlopeGaussian.pdf} 
\end{tabular}
\end{center}


\end{frame}

\begin{frame}{Le modèle linéaire généralisé mixte}

\bigskip
Notons $\mu_{it} = E(Y_{it} \mid \gamma_i)$ ($\neq E(Y_{it})$).

\bigskip
Modèle (conditionnel à $\gamma_i$, un MLG standard) :
\begin{itemize}
\item $g(E(Y_{it} \mid \gamma_i)) = g(\mu_{it}) =  x_{it}^T \beta + z_{it}^T \gamma_i$
  pour une {\color{blue}fonction de lien} (\textit{link function}) $g$.
\item $Var(Y_{it} \mid \gamma_i) = v(\mu_{it})$.
\end{itemize}

\bigskip
Deux sources d'erreur : le processus d'échantillonnage et le processus de mesure
au sein de l'échantillon lui-même.
\end{frame}

\begin{frame}{Estimation du MLGM}

Combiner l'information provenant de
\begin{itemize}
\item la distribution conditionnelle de $Y_{it}
\mid \gamma_i$ : $f(y_{it} \mid \gamma_i, \beta, \phi)$ ;
\item la distribution marginale des effets aléatoires : $f(\gamma_i \mid
\phi, \Psi)$.
\end{itemize}

\bigskip
Les $\gamma_i$ sont non observables, on s'appuie sur la {\color{blue}vraisemblance marginale} (\textit{marginal likelihood}) :
\begin{eqnarray*}
\lefteqn{L(\beta,\phi, \Psi) = f(y_1,\ldots,y_n)= \prod_{i=1}^n f(y_i) }\\
& = & \prod_{i=1}^n \int \cdots \int f(y_{i} \mid \gamma_i, \beta, \phi)f(\gamma_i \mid \phi, \Psi) d \gamma_i  \\
& = & \prod_{i=1}^n \int \cdots \int \prod_{t=1}^{n_i} f(y_{it} \mid \gamma_i,
\beta, \phi) f(\gamma_i \mid \phi, \Psi) d \gamma_i.
\end{eqnarray*}

Comme dans les MLG, utiliser la relation du modèle pour introduire l'information des données.

\end{frame}

\begin{frame}{Approximations de $L(\beta,\phi, \Psi)$}

La vraisemblance ne peut pas -- en général -- être exprimée sous forme fermée.
Des méthodes numériques et des approximations pour calculer/approximer les intégrales sont nécessaires :
\begin{itemize}
\item Intégrations de Monte-Carlo (MCMC) (approximer les intégrales par simulation)
\item {\color{blue}Quadrature de Gauss-Hermite} (\textit{Gauss-Hermite quadrature})
\item {\color{blue}Approximation de Laplace} (\textit{Laplace approximation})
\item {\color{blue}Quasi-vraisemblance pénalisée} (\textit{Penalized Quasi-Likelihood}, PQL)
\item {\color{blue}Quasi-vraisemblance marginale} (\textit{Marginal Quasi-Likelihood}, MQL)
\end{itemize}

\bigskip
Note : PQL et MQL peuvent produire des estimateurs biaisés. Il est préférable d'utiliser la
quadrature gaussienne ou l'approximation de Laplace, même si elles peuvent être coûteuses
en calcul.

\bigskip
Pour des spécifications élaborées des effets aléatoires, les seuls schémas
d'intégration efficaces semblent être ceux basés sur les MCMC.

\end{frame}

\begin{frame}{Quelques options disponibles dans R}

\begin{itemize}
\item package \texttt{nlme} : fonctions \texttt{lme} et \texttt{nlme}
  pour les modèles à effets mixtes linéaires et non linéaires. Permet des effets
  aléatoires emboîtés.
\item package \texttt{lme4} : fonction \texttt{lmer}. Permet l'approximation de Laplace (par défaut) et la
quadrature gaussienne adaptative. Peut définir des ordonnées à l'origine
aléatoires et des pentes aléatoires, ainsi que des facteurs de regroupement
emboîtés et croisés pour les effets aléatoires.
\item package \texttt{ADMB} : fonction \texttt{glmmadmb} basée sur une approximation
de Laplace avec différentiation automatique. Large gamme de familles,
large gamme de fonctions de lien, effets aléatoires simples ou multiples, incluant
des effets emboîtés et croisés.
\item package \texttt{glmmTMB}, fonction \texttt{glmmTMB}. Similaire à \texttt{glmmadmb}, mais avec certaines extensions différentes.
\end{itemize}

\end{frame}


\begin{frame}{Interprétation des effets fixes}
Un MLGM est une {\color{blue}analyse spécifique au sujet} (\textit{subject-specific analysis}) : les coefficients
de régression s'appliquent à chaque individu, mais pas nécessairement à la population (par
opposition à une {\color{blue}analyse marginale} (\textit{marginal analysis}), par ex.\ les GEE).

\bigskip
On a :
$$\mu_{it} = E(Y_{it}\mid \gamma_i) = g^{-1}(x_{it}^T \beta + z_{it}^T \gamma_i),$$
de sorte que la moyenne conditionnelle varie avec les prédicteurs selon
$\beta$ avec en plus une contribution spécifique au sujet due à $\gamma_i$.

\end{frame}

\begin{frame}
La moyenne populationnelle $E(Y_{it})$ est définie par
$$E(Y_{it}) = E_{\gamma_i}(E(Y_{it}\mid \gamma_i)) = \int g^{-1}(x_{it}^T \beta + z_{it}^T
\gamma_i) dF(\gamma_i),$$
qui est différente de $g^{-1}(x_{it}^T \beta)$, sauf dans le
cas gaussien avec lien identité.

\bigskip
Dans certains cas (par ex.\ le modèle probit), on a
$$E(Y_{it}) = E_{\gamma_i}(E(Y_{it}\mid \gamma_i)) = g^{-1}(x_{it}^T \beta^*),$$
avec $\mid \beta^* \mid \leq  \mid \beta \mid$.
\end{frame}

\begin{frame}
\begin{center}
\begin{tabular}{cc} 
Ordonnée à l'origine aléatoire & Pente aléatoire \\
\includegraphics[width = 5.5cm]{figures/RandomIntercept.pdf}  & \includegraphics[width =5.5cm]{figures/RandomSlope.pdf} 
\end{tabular}
\end{center}

\end{frame}

\begin{frame}{Prédiction des effets aléatoires}
Pour les modèles mixtes linéaires, il existe une théorie pour définir les {\color{blue}meilleurs prédicteurs} (\textit{best predictors}, BP) minimisant
l'EQM, $\hat{\gamma}_i = E(\gamma_i \mid y_i)$, les {\color{blue}meilleurs prédicteurs linéaires} (\textit{best linear
predictors}, BLP) et les {\color{blue}meilleurs prédicteurs linéaires sans biais} (\textit{best linear unbiased predictors}, BLUP). Ils reviennent tous
au même.

\bigskip
Dans les logiciels pour MLGM, ce sont généralement les {\color{blue}modes conditionnels} (\textit{conditional modes}) qui sont implémentés, c'est-à-dire les
valeurs maximisant la densité de $\gamma_i \mid y_1, \ldots, y_n$. Ce ne sont pas
des BLUP, car ils ne sont pas linéaires, ils pourraient ne pas être sans biais et on ne sait pas
en quel sens ils sont « meilleurs ».

\bigskip
Les meilleurs prédicteurs seraient ceux satisfaisant
$\hat{\gamma}_i = E(\gamma_i \mid y_i)$. C'est encore un domaine de recherche ouvert (en
particulier pour définir les propriétés de $\hat{\gamma}_i$).
\end{frame}

\begin{frame}
En fait, la question à poser est quelle est la prédiction pour :
\begin{itemize}
\item un nouveau point de données dans un groupe existant : utiliser $\hat{\mu}_{it} = g^{-1}(x_{it}^T \hat{\beta} + z_{it}^T \hat{\gamma}_i)$ ;
\item un nouveau point de données pour un nouveau groupe : utiliser $\hat{\mu}_i = g^{-1}(x_{it}^T
\hat{\beta})$.
\end{itemize}
\end{frame}

\subsection{Jeu de données sur la pseudonégligence}
\begin{frame}{Analyse des données sur la pseudonégligence}

Les effets de l'hémiespace sur une tâche de bissection tactile de ligne sont mesurés. Les données proviennent d'un mémoire de diplôme (Sabina Catalano, FAPSE, 2003). On dispose d'observations sur 10 sujets mesurés 10 fois pour chacune des 13 distances (de 0 à 120, 60 étant le centre de la ligne). Le sujet se voit présenter (dans un ordre aléatoire) une position sur une ligne et doit décider si elle se situe à gauche ou à droite du centre. Le nombre de réponses « gauche » (pour chaque distance) est enregistré.

Dans cette expérience, la question scientifique est de déterminer s'il existe un biais systématique vers la gauche par rapport au centre de la ligne (comme suggéré par certains psychologues).

\bigskip
Le {\color{blue}point d'égalité subjective} (\textit{equality subjective point}, ESP), c'est-à-dire la distance
pour laquelle la probabilité de « gauche » est 1/2, est le paramètre d'intérêt.
\end{frame}

\begin{frame}
\begin{center}
\includegraphics[width = 7.5cm]{figures/hemiID.pdf}  
\end{center}

\end{frame}

\begin{frame}{Moyenne sur les individus}

\begin{center}
\includegraphics[width = 7cm]{figures/hemi.pdf}  
\end{center}
\end{frame}

\begin{frame}
On considère un modèle MLGM binomial avec ordonnée à l'origine aléatoire :

$$\mbox{logit}(p_i) = \log\Big( \frac{p_i}{1-p_i}\Big) = \beta_0 + \beta_1 \texttt{distance} + \gamma_i,$$
avec $p_i = P(\mbox{« gauche » pour l'individu }i)$ et $\gamma_i \sim {\mathcal N}(0,\sigma_{\gamma}^2)$ est une ordonnée à l'origine aléatoire.

\bigskip
\bigskip
\bigskip
Pour chaque individu, le ESP est la distance telle que
$$ \beta_0 + \beta_1 \texttt{distance} + \gamma_i =0,$$

c'est-à-dire
$$ ESP_i = - \frac{\beta_0 + \gamma_i}{\beta_1}.$$
\end{frame}


\begin{frame}[fragile]
{\tiny
\begin{verbatim}
Generalized linear mixed model fit by maximum likelihood (Laplace Approximation) ['glmerMod']
 Family: binomial  ( logit )
Formula: cbind(left, 10 - left) ~ distance + (1 | sujet)
   Data: hemi

     AIC      BIC   logLik deviance df.resid 
   371.4    380.0   -182.7    365.4      127 

Scaled residuals: 
    Min      1Q  Median      3Q     Max 
-6.6927 -0.6634 -0.1685  0.6606  3.5747 

Random effects:
 Groups Name        Variance Std.Dev.
 sujet  (Intercept) 0.2469   0.4969  
Number of obs: 130, groups:  sujet, 10

Fixed effects:
             Estimate Std. Error z value Pr(>|z|)    
(Intercept)  3.804820   0.264566   14.38   <2e-16 ***
distance        -0.063078   0.003267  -19.31   <2e-16 ***
---
Signif. codes:  0 '***' 0.001 '**' 0.01 '*' 0.05 '.' 0.1 ' ' 1

Correlation of Fixed Effects:
     (Intr)
distance -0.744
\end{verbatim}
}
\end{frame}

\begin{frame}
Avec prudence :
\begin{center}
\includegraphics[width = 7cm]{figures/hemires.pdf}  
\end{center}
\end{frame}


\begin{frame}[fragile]
Les effets aléatoires en valent-ils la peine ? Comparaison avec un ajustement MLG (qui est emboîté dans un modèle MLGM) :

{\tiny
\begin{verbatim}
Call: glm(formula = cbind(gauche, 10 - gauche) ~ item, family = binomial, 
    data = hemi)

Deviance Residuals: 
    Min       1Q   Median       3Q      Max  
-3.1932  -0.9789  -0.2511   0.9206   3.6980  

Coefficients:
             Estimate Std. Error z value Pr(>|z|)    
(Intercept)  3.650109   0.202573   18.02   <2e-16 ***
item        -0.060521   0.003093  -19.57   <2e-16 ***
---
Signif. codes:  0 '***' 0.001 '**' 0.01 '*' 0.05 '.' 0.1 ' ' 1

(Dispersion parameter for binomial family taken to be 1)

    Null deviance: 988.40  on 129  degrees of freedom
Residual deviance: 198.04  on 128  degrees of freedom
AIC: 390.75

Number of Fisher Scoring iterations: 5
\end{verbatim}
}

\bigskip
Dans le MLGM : AIC=371.4 et la déviance=365.4.
\end{frame}




\begin{frame}[fragile]{Prédiction du ESP}

{\tiny
\begin{verbatim}
-(fixef(hemi.glmm)[1] +ranef(hemi.glmm)$sujet)/fixef(hemi.glmm)[2]
        (Intercept)
Sujet1     55.97849
Sujet10    51.06305
Sujet2     58.43176
Sujet3     59.24924
Sujet4     60.88410
Sujet5     77.32295
Sujet6     51.88350
Sujet7     63.33696
Sujet8     61.70159
Sujet9     63.33696
\end{verbatim}
}

\vspace*{-2cm}
\begin{center}
\includegraphics[width = 6cm]{figures/ESP.pdf}  
\end{center}
\end{frame}


\begin{frame}[fragile]{Avec pente aléatoire ?}

{\tiny
\begin{verbatim}
Generalized linear mixed model fit by maximum likelihood (Laplace Approximation) ['glmerMod']
 Family: binomial  ( logit )
Formula: cbind(left, 10 - left) ~ item + (1 + distance | sujet)
   Data: hemi

     AIC      BIC   logLik deviance df.resid 
   368.4    382.8   -179.2    358.4      125 

Scaled residuals: 
    Min      1Q  Median      3Q     Max 
-3.8813 -0.5607 -0.0854  0.5181  2.3284 

Random effects:
 Groups Name        Variance  Std.Dev. Corr 
 sujet  (Intercept) 1.2714312 1.12758       
        distance    0.0001859 0.01363  -0.87
Number of obs: 130, groups:  sujet, 10

Fixed effects:
            Estimate Std. Error z value Pr(>|z|)    
(Intercept)  4.02204    0.43108    9.33   <2e-16 ***
distance    -0.06629    0.00565  -11.73   <2e-16 ***
---
Signif. codes:  0 '***' 0.001 '**' 0.01 '*' 0.05 '.' 0.1 ' ' 1

Correlation of Fixed Effects:
     (Intr)
item -0.885
\end{verbatim}
}
\end{frame}

\begin{frame}{Remerciements}
Ces diapositives ont été inspirées par le cours GLAM de Eva Cantoni de l'Université de Genève.
\end{frame}

\begin{frame}[allowframebreaks]{Références du chapitre}
\tiny
\printbibliography[heading=none,section=\therefsection]
\end{frame}

\end{refsection}

\end{document}
